\section{Background and Foundations}\label{sec:background}

The reference pipeline proposed in this article assembles several mature technologies into a new lifecycle role, and its design choices are intelligible only against the formal-methods machinery they reuse. This section establishes that machinery. It proceeds from the question of \emph{what a specification is} and how design by contract renders one machine-checkable in the Java Modeling Language (JML); through \emph{how such a contract is discharged} by deductive verification, extended static checking (ESC), OpenJML, and an underlying Satisfiability Modulo Theories (SMT) solver, together with the decidability and scalability limits that constrain what can be proved; to \emph{how a contract can be obtained cheaply} through specification inference, of which the author's JML-Inferrer instrument is the concrete realisation carried forward into this work; and finally to the \emph{generators} of code---large language models (LLMs) and the agents built on them---whose output the pipeline is designed to gate. Throughout, the intent is didactic: each foundation is introduced with the precise terminology used in the remainder of the paper, so that the contributions of \cref{sec:pipeline} can be stated without re-deriving their substrate.

\subsection{Specifications and design by contract}\label{subsec:bg-dbc}

A \emph{specification} is a statement of intended behaviour that is independent of any particular implementation. In the formal-methods tradition this statement is a predicate, or a pair of predicates, over program states. The axiomatic basis is Hoare logic, in which a Hoare triple $\{P\}\,S\,\{Q\}$ asserts that if the precondition $P$ holds before the statement $S$ executes and $S$ terminates, then the postcondition $Q$ holds afterward~\cite{hoare1969axiomatic}. Dijkstra's predicate-transformer semantics recasts this as the weakest precondition $\mathrm{wp}(S, Q)$, the weakest predicate on the initial state guaranteeing that $S$ terminates and establishes $Q$ (total correctness; the partial-correctness reading of the triple above corresponds to the weakest \emph{liberal} precondition), giving a calculus for reasoning about correctness mechanically rather than by inspection~\cite{dijkstra1976discipline,dijkstra1975guarded}. The refinement calculus extends this into a methodology in which an abstract specification is transformed, step by provably correct step, into executable code~\cite{back1998refinement,morgan1994programming}; that tradition is the intellectual ancestor of the contract-to-code \emph{synthesis} discussed later, with the crucial difference that the pipeline of this article delegates the transformation to a fallible agent and recovers soundness by verifying the result rather than by trusting the derivation.

Design by Contract (DbC), introduced by Meyer for Eiffel, operationalizes these ideas for working object-oriented software~\cite{meyer1992dbc,meyer1997oosc}. A method is governed by a contract comprising a precondition that the caller must establish, a postcondition that the method guarantees in return, and class invariants that every public method must preserve. The arrangement assigns \emph{obligations} and \emph{benefits} between caller and callee: the caller is obliged to satisfy the precondition and benefits from the postcondition; the callee benefits from assuming the precondition and is obliged to deliver the postcondition. This caller/callee duality is foundational to the present work, because it is exactly what makes contracts \emph{compose}: a change to a callee's contract recomputes the obligations of every caller, the property the thesis programme exploits for version-compatibility analysis (\cref{subsec:bg-inference}).

JML is the behavioural interface specification language for Java in which the pipeline expresses its contracts~\cite{leavens2006jml}. JML annotations are written in specially marked comments---so they are ignored by the Java compiler---and use a superset of Java's expression syntax. A \texttt{requires} clause states a precondition and an \texttt{ensures} clause a postcondition; \texttt{invariant} declares a class invariant; an \texttt{assignable} (frame) clause bounds the locations a method may modify, which is essential for sound modular reasoning because it tells the verifier what a call \emph{cannot} change. Postconditions may refer to the pre-state value of an expression through \verb|\old(e)| and to the returned value through \verb|\result|, and may quantify over ranges with \verb|\forall| and \verb|\exists|. \Cref{lst:jml-example} illustrates the vocabulary on a small bounded-search method; the contract states that the array is non-null, that a non-negative result indexes an element equal to the key, and that a negative result certifies the key is absent.

\begin{lstlisting}[caption={A JML method contract using the core clauses the pipeline manipulates.},label={lst:jml-example},captionpos=b]
//@ requires a != null;
//@ ensures \result >= 0 ==> (\result < a.length && a[\result] == key);
//@ ensures \result <  0 ==> (\forall int k; 0 <= k && k < a.length; a[k] != key);
//@ assignable \nothing;
int indexOf(int[] a, int key) { ... }
\end{lstlisting}

Two properties of this representation matter for the pipeline. First, a JML contract is \emph{machine-checkable}: unlike a natural-language docstring, it has a formal semantics against which an implementation can be mechanically judged. Second, it is \emph{partial}: a contract constrains behaviour only to the extent its clauses are written, so an under-specified contract may be satisfied by implementations the author never intended---a weakness this article treats as a first-class hazard under the headings of vacuity and spec weakening (\cref{sec:pipeline}). The closest contemporary relative of JML is Spec\#, which embeds contracts in C\# and discharges them with a verifying compiler~\cite{barnett2005specsharp}; the design considerations transfer, but JML's tooling maturity for Java is what makes it the substrate here.

\subsection{Deductive verification, ESC, OpenJML, and SMT}\label{subsec:bg-verification}

Possessing a machine-checkable contract is useful only if an implementation can be \emph{checked} against it across all inputs. Deductive program verification does this by reducing the question ``does this code satisfy this contract?'' to a set of logical formulas---\emph{verification conditions} (VCs)---whose validity entails correctness. The VCs are generated by a weakest-precondition computation over the method body~\cite{dijkstra1976discipline}: the verifier propagates the postcondition backward through each statement, accumulating the obligations that must hold, and emits a formula asserting that the precondition implies them. Loops break this backward walk because their effect over an unknown number of iterations cannot be enumerated; the engineer (or an inference tool) must supply a \emph{loop invariant}, a predicate true on entry, preserved by each iteration, and strong enough that together with the loop's exit condition it entails the obligation the loop must establish, which the verifier uses in place of unrolling. The need for loop invariants is the single largest source of the annotation burden that motivates specification inference.

Extended Static Checking is the pragmatic engineering of this idea pioneered by ESC/Java~\cite{flanagan2002escjava}. ESC trades completeness for automation and modularity: it analyses each method in isolation, assuming the contracts of the methods it calls rather than inlining them, and it discharges the resulting VCs automatically without interactive proof. OpenJML is the modern realisation of ESC for Java built on JML~\cite{cok2011openjml,cok2014openjml}; it parses JML annotations, generates VCs from method bodies under their contracts, and hands those VCs to an automated theorem prover. That prover, in the configuration used by this article's instrument, is the Z3 SMT solver~\cite{demoura2008z3}, which determines the satisfiability of formulas over combined background theories---linear integer and real arithmetic, arrays, bit-vectors, and uninterpreted functions---that match the constructs arising from Java programs; the quantifier-free fragments of these theories are decidable, whereas the quantified formulas arising from JML contracts are handled by heuristic instantiation and may return \texttt{unknown} or fail to terminate. cvc5 is a comparable alternative~\cite{barbosa2022cvc5}. OpenJML asks Z3 whether the negation of each VC is satisfiable: an \texttt{unsat} answer certifies the VC and hence the obligation, while a satisfiable answer yields a model that OpenJML reports as a \emph{counterexample}---a concrete input witnessing the violation. This counterexample is not a diagnostic convenience but a structural asset of the pipeline: it is a precise, machine-generated witness of failure that the synthesis loop feeds back to the code generator---typically a far more informative signal than ``a test failed,'' though under modular checking such a witness can be spurious (for instance, reflecting a loop invariant too weak to carry the proof rather than a reachable defect).

It must be stated precisely what such a proof does and does not establish, because overclaiming here would undermine the entire argument. A successful ESC run proves that the implementation satisfies its contract for \emph{all} inputs admitted by the precondition---a guarantee no finite test suite can give. It does \emph{not} prove the contract captures what the developer wanted; that is the province of the intent layer. Equally important are the limits. ESC verifies modularly under assumed callee contracts, so its guarantee is only as strong as those contracts. Program verification is undecidable in general, and even the decidable fragments Z3 targets carry combinatorial cost: in practice the solver times out on nonlinear (multiplicative) arithmetic, on rich reasoning over the array theory, and on contracts with deeply nested quantifiers. The thesis programme's own corpus exhibits exactly these failure modes, and the pipeline's design accommodates them through \emph{tiered assurance} (\cref{sec:pipeline}) rather than pretending verification always terminates. The author's instrument additionally relies on a custom OpenJML fork that adds recursive function definitions (\verb|define-fun-rec|) for the JML aggregation operators \verb|\sum|, \verb|\product|, and \verb|\num_of|, which the stock toolchain does not discharge; this fork is reused as a trusted component in the proof-of-concept of \cref{sec:poc}. Comparable deductive ecosystems---Dafny's auto-active verifier~\cite{leino2010dafny}, the KeY interactive prover for Java~\cite{ahrendt2016key}, Frama-C for C~\cite{kirchner2015framac}, and Verus for Rust~\cite{lattuada2023verus}---share the same VC-to-SMT architecture and the same fundamental limits, confirming that the constraints discussed here are intrinsic rather than artefacts of one tool.

\subsection{Specification inference and the JML-Inferrer}\label{subsec:bg-inference}

The chief obstacle to deploying contract-based verification is the cost of \emph{authoring} the contracts. Specification inference attacks this bootstrap cost by deriving candidate specifications automatically. Two traditions exist. Dynamic inference, exemplified by Daikon, observes program executions and reports invariants that held over the traces as \emph{likely} program properties~\cite{ernst2007daikon}; the inferred properties are only as representative as the test inputs and require subsequent validation, and integrating them with static checkers is itself a research thread~\cite{nimmer2001static}. Static inference instead reasons about the code without running it: Houdini conjectures a large pool of candidate annotations and iteratively retracts those a checker refutes~\cite{flanagan2001houdini}; abstract interpretation computes sound over-approximations of reachable states~\cite{cousot1977abstract}; and dedicated invariant-learning frameworks such as ICE~\cite{garg2014ice} and neural approaches like Code2Inv~\cite{si2018learning} target the loop-invariant problem specifically. A separate strand mines specifications from execution traces as API-protocol automata~\cite{ammons2002mining}.

The instrument carried into this work, the \emph{JML-Inferrer}, occupies the static end of this space. It is a Java~21 tool built on the JavaParser library that analyses a method's abstract syntax tree (AST) and emits JML annotations---preconditions, postconditions, loop invariants, class invariants, and \texttt{assignable} clauses---by heuristic pattern matching over AST structure rather than by a complete program analysis. A null-check or array access on a parameter induces a non-nullness precondition; recognised loop shapes (counting up or down, accumulating, linear search, running maximum or minimum) induce matching invariants and result postconditions; numeric guards induce range and overflow preconditions. The heuristic basis is a deliberate trade: the inferrer favors recall of plausible, syntactically valid JML over the soundness guarantees of abstract interpretation, on the premise that an automated checker (\cref{subsec:bg-verification}) will reject any inferred clause that does not hold. This instrument is not introduced here; it is the validated artefact of the author's prior programme, reused as the brownfield contract-bootstrapping component of the proposed pipeline.

Four results from that programme are load-bearing for the present article and are cited as the author's own prior work. Article~1 established that automatically inferred JML specifications materially improve LLM-based unit-test generation, the first evidence that inferred contracts have downstream value for AI-assisted development~\cite{edmonds_article1}. Article~2 showed that JML specifications can be embedded into compiled Java artefacts and into OpenAPI service descriptions with full roundtrip fidelity, so a contract can travel with the code it governs~\cite{edmonds_article2}. Article~3 demonstrated that heuristic inference admits \emph{compositional refinement} across method boundaries---a second weakest-precondition pass over the call graph recovers preconditions a single-method pass misses, contributing on the order of $18{,}854$ additional well-formed preconditions on the experimental corpus~\cite{edmonds_article3}. Because contracts compose along the caller/callee obligation structure of DbC (\cref{subsec:bg-dbc}), this same mechanism lets a callee-spec change recompute caller obligations \emph{statically}, which is the technical bridge from inference to the version-compatibility capability of \cref{sec:pipeline}. Article~4 integrated the inference pipeline into continuous-integration systems with diff-only analysis, content-hash caching, and fail-open semantics, establishing that the toolchain runs at the cadence of ordinary development~\cite{edmonds_article4}. Together these results supply the reused, trusted components on which the proof-of-concept of \cref{sec:poc} is built, and they ground the present article's claim that the missing piece is not the verification machinery but its \emph{lifecycle organisation}.

\subsection{LLMs and agents as code generators}\label{subsec:bg-llm}

The fourth foundation is the code generator the pipeline gates. LLMs trained on source code can map a natural-language description to plausible code: Codex demonstrated function-level generation from docstrings and introduced the \emph{pass@k} functional-correctness metric~\cite{chen2021codex}, with MBPP and AlphaCode extending the evidence to entry-level and competition-level tasks respectively~\cite{austin2021mbpp,li2022alphacode}. These systems inverted the economics of classical synthesis---natural-language intent is cheap and abundant---but estimate correctness against finite, model-adjacent test suites rather than proving it against a contract. That estimate is fragile: benchmark suites are often too weak to expose plausible-but-wrong solutions, and strengthening them lowers reported pass rates~\cite{liu2023evalplus}; and generated code is prone to confident fabrication of non-existent APIs and subtly inconsistent logic that survives superficial review~\cite{zhang2025hallucination}.

The frontier has since moved from isolated functions to \emph{agents} that operate over whole repositories. SWE-bench reframed evaluation around resolving real GitHub issues with project test suites as the oracle~\cite{jimenez2024swebench}; SWE-agent showed that a tailored agent--computer interface markedly improves repository navigation and editing~\cite{yang2024sweagent}; and AutoCodeRover paired LLM reasoning with structure-aware search and fault localization to localize and patch issues~\cite{zhang2024autocoderover}. A complementary line equips agents with iterative self-correction---Self-Refine~\cite{madaan2023selfrefine}, Reflexion~\cite{shinn2023reflexion}, execution-trace-driven self-debugging~\cite{chen2024selfdebug}, and the test-driven flow of AlphaCodium~\cite{ridnik2024alphacodium}---while multi-agent frameworks such as MetaGPT~\cite{hong2024metagpt}, ChatDev~\cite{qian2024chatdev}, and AutoGen~\cite{wu2024autogen} coordinate role-specialized agents over a development workflow, as surveyed recently for software engineering~\cite{liu2024llmagentse}. These advances are real, but they preserve two structural defects: the specification of intent remains natural-language issue text, and the acceptance gate remains an executable test suite---frequently the project's own, which the agent may also consult or edit, so the model effectively grades itself. This circularity is acute in oracle generation, where the same model authors both the code and the tests meant to judge it~\cite{liu2026oracle}.

These defects are being institutionalized rather than resolved. Emerging \emph{spec-driven} lifecycles---AWS's AI-Driven Development Lifecycle~\cite{awsaidlc}, GitHub's Spec Kit~\cite{githubspeckit}, and the Kiro IDE~\cite{awskiro}, all reflecting the broader repositioning of LLMs within the software lifecycle~\cite{ozkaya2023llmse}---reorganise development around a specification, but that specification is natural language, the loop from specification to code is unverifiable, and verification, where present, remains self-grading. The contribution of this article, developed in \cref{sec:pipeline}, is to interpose a machine-checkable formal contract as the pivot of these lifecycles: a single artefact that serves simultaneously as the input intent an agent builds against and the oracle its output is verified against by OpenJML over all inputs, with an independent-authority constraint and provenance- and assurance-tagged obligations ensuring the verifier never grades its own author. The foundations assembled here---contracts and their composition, ESC and its limits, inference and its prior validation, and the agentic generators and their circularity---are exactly the pieces that pipeline re-couples.
